% 2026-11-16 站主問：「Constitutional theory is dead; long live constitutional theory. 隱喻來源」「對於我們講違憲審查正當性的理解有什麼幫助」。
% Law 800--801 頁主對話自文字層讀過；Fallon 21--22 頁自 readings/W04/text/ 讀過。
% 套語之首次使用年份：通行的 1422 年（查理七世繼位）一說，英文維基該節標示無出處；王室葬禮傳令官宣告之記載未取得一手文獻，正文不寫年份。
\chapter{結語之套語與違憲審查之正當性}
\label{ch:50_theory}

\begin{qaquestion}
Law\ix{persons}{law@Law, David S.} 全文最後一句是「Constitutional theory is dead; long live constitutional theory」（801 頁）。
\begin{enumerate}
\tightlist
\item 這句話套用的是什麼成語？
\item Law 用它宣告什麼死去、什麼延續？
\item 對理解違憲審查正當性的討論有何幫助？
\end{enumerate}
\end{qaquestion}

\section{套語的來源}
\label{sec:50_theory-1}

句式套用法國王位繼承的宣告「Le roi est mort, vive le roi」〔國王駕崩，國王萬歲〕。前一句的國王是剛去世的前王，後一句的國王是即位的新王；前王一死，王位即由繼承人取得，其間沒有空位。法諺 le mort saisit le vif〔死者使生者取得〕表達的是同一種承繼：權利於死亡的同一時點移轉於繼承人。\footcite{wiki-roi}

\section{Law 的用法}
\label{sec:50_theory-2}

這句話之前，Law 寫道：「But it is time for the wheel to turn. The tired tropes of the countermajoritarian dilemma have occupied too much attention for too long」〔然而風水該輪流轉了。抗多數難題那些陳腐的套話，佔據注意力太多、太久〕。\footcite[801]{law2009}800 頁註 242 引 Friedman 所稱對抗多數難題的「academic obsession」〔學術執念〕。\footcite[800 n.242]{law2009}

死去的是以抗多數難題\ix{terms}{02@\zhlidxcat{司法權之基礎}!抗多數難題}為中心問題的憲法理論。Law 在 727 頁引 Bickel\ix{persons}{bickel@Bickel, Alexander M.}：最高法院宣告民選官員的行為違憲時，「exercises control, not in behalf of the prevailing majority, but against it」〔所行使的控制，不是代表當時的多數，而是反對它〕；此後憲法理論即由法官何時得推翻多數的問題所支配。\footcite[727]{law2009}繼位的是 800 至 801 頁所列的研究問題：人民為何服從法院、法院發揮何種功能而使誰受益、政府為何遵從憲法法院的裁判、法院能做什麼以擴張其權力、人民與政府何以不只容忍而且鼓勵司法權的擴張。\footcite[800--801]{law2009}依套語的結構，憲法理論這一學科延續不斷，換掉的是它的中心問題。

\section{對違憲審查正當性討論的幫助}
\label{sec:50_theory-3}

\begin{enumerate}
\tightlist
\item
  W04 Fallon\ix{persons}{fallon@Fallon, Richard H., Jr.} 區分三種正當性概念\ix{terms}{02@\zhlidxcat{司法權之基礎}!正當性}。社會學意義的正當性取決於公眾對政府、機關或裁判的實際態度與反應，與他們應當如何想無關；道德意義的正當性涉及人們不論他人如何想都應當持有的態度；法律意義的正當性則是法律概念。Fallon 指出，一個政權或一則裁判可以廣受贊同而在道德上錯誤，納粹政權即可能享有廣泛的社會學正當性。\footcite[21--22]{fallon2018}Law 列出的新問題，人民為何服從、政府為何遵從，都是社會學意義的正當性問題，也可以說是權力的問題。Law 在 779 頁批評把正當性當成一座可以累積也可以耗盡的「reservoir」〔蓄水池〕的通說，793 頁說這套說法迴避了服從的動機。\footcite[779, 793]{law2009}他所宣告死去的，是混用三種概念的正當性討論；他的協調理論回答的是其中一種。
\item
  Bush v. Gore\ix{cases}{02@美國!01@聯邦最高法院!2000Bush v. Gore（2000）@Bush v. Gore（2000）} 一案在兩篇中的評價可資對照。Fallon 記載，許多贊同 Roe 的人主張 Bush v. Gore 不具正當性，也有許多論者預料兩案會使最高法院失去更一般意義的正當性。\footcite[20--21]{fallon2018}Law 則說 Bush v. Gore 在強烈批評與廣泛反對之下協調了和平的政權交替，使人普遍相信法院能夠再做一次，因而是法院權力的增加。\footcite[782]{law2009}兩人對同一案件的判斷不同，所用的正當性概念不同：批評者談的是法律與道德意義的正當性，Law 量的是服從，即社會學意義的正當性。
\item
  抗多數難題預設法院與多數對立。Law 在 731、786 頁把委託人改定為人民、把政府與立法多數定為代理人，法院監督代理人，於是成為人民主權的工具。\footcite[731, 786]{law2009}抗多數難題在這個框架下消解，W06 Waldron\ix{persons}{waldron@Waldron, Jeremy} 提出的規範問題卻未獲回答：公民對權利內容存有合理歧見時，由法院而非立法機關作最後決定，是否正當。\footcite{waldron2006}第 1 部第四節第 3 項已指出，人人預期他人遵守的錯誤判決，同樣是穩定的協調均衡。Law 換掉了討論正當性時的提問方式，道德意義的正當性問題仍須另行回答。
\item
  課堂上討論某一號解釋或判決有無正當性時，可以先問所指的是哪一種：該裁判是否被各機關遵從（社會學意義），是否有憲法與法律上的依據（法律意義），是否應當被遵從（道德意義）。Law 的理論只能回答第一種，Fallon 的三分法可以用來檢驗任何一篇文獻在回答哪一種。
\end{enumerate}

\section{課前可提的問題}
\label{sec:50_theory-4}

\begin{enumerate}
\tightlist
\item Law 宣告以抗多數難題為中心的憲法理論已死。他的協調理論回答的是社會學意義的正當性，道德意義的正當性問題由誰繼承？
\item 依 Fallon 的三分法，Law 782 頁把 Bush v. Gore 計為法院權力的增加，與批評者認為該案不具正當性，兩者是否真有衝突？
\end{enumerate}
